% Lookback choices in neural-network volatility forecasting studies.
% Drop into main.tex with % Lookback choices in neural-network volatility forecasting studies.
% Drop into main.tex with % Lookback choices in neural-network volatility forecasting studies.
% Drop into main.tex with % Lookback choices in neural-network volatility forecasting studies.
% Drop into main.tex with \input{lookback_table}. Requires booktabs.
% Compiled 29 Oct 2026 from the papers' full texts; checked again 11 Nov 2026.
\begin{table}[!htbp]
\centering
\caption{Lookback lengths in LSTM volatility forecasting studies}
\label{tab:lookbacks}
\footnotesize
\setlength{\tabcolsep}{4pt}
\begin{tabular}{@{}p{3.6cm}p{3.9cm}p{2.0cm}p{5.6cm}@{}}
\toprule
Study & Data & Lookback (days) & How the lookback was chosen \\
\midrule
\textcite{petrozziello-2022} & 28 DJIA stocks and SPY; NASDAQ-100 stocks & 20 & Grid search over 5--100 on the first year of data; a single value reported \\
\textcite{chen-2022} & CSI 300 and S\&P 500 index futures & 20 & ``After some initial tests''; tests not reported \\
\textcite{rosenbaum-2022} & 862 US stocks, pooled & 22 & Stated, ``to make it comparable with the HAR model'' \\
\textcite{zhang-2023} & US stocks, intraday RV, pooled & 21 & Stated; the same 21-day window for every model \\
\textcite{roszyk-2024} & S\&P 500 index & 22 & Stated, ``the approximate number of trading days in a month'' \\
\textcite{perekhodko-2025} & S\&P 500 index & 21 & Stated; 5 and 42 checked afterwards \\
\textcite{souto-2023} & S\&P 500 index & 21, 63 or 252 & Random search over 21--252 on a validation split; differs by model \\
\textcite{tang-2024} & 213 cryptocurrencies, pooled & 7 or 30 & Cross-validation within the training sample \\
\textcite{rodikov-2022} & Stocks, S\&P 500, Bitcoin, Ether & 1--50 & Validation search; selected values not reported \\
\textcite{bucci-2020} & S\&P 500, monthly RV & Not reported & Information criteria for lags of RV in the predictor set; sequence length of the recurrent networks not stated \\
\midrule
This paper & 27 DJIA stocks & 6--63 & AIC on all data before each annual re-selection date, separately for each stock \\
\bottomrule
\end{tabular}

\vspace{0.4em}
\begin{minipage}{0.97\textwidth}
\footnotesize
\textit{Notes:} Lookback is the number of past daily observations in the LSTM input sequence. Where a study reports several values, the table gives those used in its main results.
\end{minipage}
\end{table}
. Requires booktabs.
% Compiled 29 Oct 2026 from the papers' full texts; checked again 11 Nov 2026.
\begin{table}[!htbp]
\centering
\caption{Lookback lengths in LSTM volatility forecasting studies}
\label{tab:lookbacks}
\footnotesize
\setlength{\tabcolsep}{4pt}
\begin{tabular}{@{}p{3.6cm}p{3.9cm}p{2.0cm}p{5.6cm}@{}}
\toprule
Study & Data & Lookback (days) & How the lookback was chosen \\
\midrule
\textcite{petrozziello-2022} & 28 DJIA stocks and SPY; NASDAQ-100 stocks & 20 & Grid search over 5--100 on the first year of data; a single value reported \\
\textcite{chen-2022} & CSI 300 and S\&P 500 index futures & 20 & ``After some initial tests''; tests not reported \\
\textcite{rosenbaum-2022} & 862 US stocks, pooled & 22 & Stated, ``to make it comparable with the HAR model'' \\
\textcite{zhang-2023} & US stocks, intraday RV, pooled & 21 & Stated; the same 21-day window for every model \\
\textcite{roszyk-2024} & S\&P 500 index & 22 & Stated, ``the approximate number of trading days in a month'' \\
\textcite{perekhodko-2025} & S\&P 500 index & 21 & Stated; 5 and 42 checked afterwards \\
\textcite{souto-2023} & S\&P 500 index & 21, 63 or 252 & Random search over 21--252 on a validation split; differs by model \\
\textcite{tang-2024} & 213 cryptocurrencies, pooled & 7 or 30 & Cross-validation within the training sample \\
\textcite{rodikov-2022} & Stocks, S\&P 500, Bitcoin, Ether & 1--50 & Validation search; selected values not reported \\
\textcite{bucci-2020} & S\&P 500, monthly RV & Not reported & Information criteria for lags of RV in the predictor set; sequence length of the recurrent networks not stated \\
\midrule
This paper & 27 DJIA stocks & 6--63 & AIC on all data before each annual re-selection date, separately for each stock \\
\bottomrule
\end{tabular}

\vspace{0.4em}
\begin{minipage}{0.97\textwidth}
\footnotesize
\textit{Notes:} Lookback is the number of past daily observations in the LSTM input sequence. Where a study reports several values, the table gives those used in its main results.
\end{minipage}
\end{table}
. Requires booktabs.
% Compiled 29 Oct 2026 from the papers' full texts; checked again 11 Nov 2026.
\begin{table}[!htbp]
\centering
\caption{Lookback lengths in LSTM volatility forecasting studies}
\label{tab:lookbacks}
\footnotesize
\setlength{\tabcolsep}{4pt}
\begin{tabular}{@{}p{3.6cm}p{3.9cm}p{2.0cm}p{5.6cm}@{}}
\toprule
Study & Data & Lookback (days) & How the lookback was chosen \\
\midrule
\textcite{petrozziello-2022} & 28 DJIA stocks and SPY; NASDAQ-100 stocks & 20 & Grid search over 5--100 on the first year of data; a single value reported \\
\textcite{chen-2022} & CSI 300 and S\&P 500 index futures & 20 & ``After some initial tests''; tests not reported \\
\textcite{rosenbaum-2022} & 862 US stocks, pooled & 22 & Stated, ``to make it comparable with the HAR model'' \\
\textcite{zhang-2023} & US stocks, intraday RV, pooled & 21 & Stated; the same 21-day window for every model \\
\textcite{roszyk-2024} & S\&P 500 index & 22 & Stated, ``the approximate number of trading days in a month'' \\
\textcite{perekhodko-2025} & S\&P 500 index & 21 & Stated; 5 and 42 checked afterwards \\
\textcite{souto-2023} & S\&P 500 index & 21, 63 or 252 & Random search over 21--252 on a validation split; differs by model \\
\textcite{tang-2024} & 213 cryptocurrencies, pooled & 7 or 30 & Cross-validation within the training sample \\
\textcite{rodikov-2022} & Stocks, S\&P 500, Bitcoin, Ether & 1--50 & Validation search; selected values not reported \\
\textcite{bucci-2020} & S\&P 500, monthly RV & Not reported & Information criteria for lags of RV in the predictor set; sequence length of the recurrent networks not stated \\
\midrule
This paper & 27 DJIA stocks & 6--63 & AIC on all data before each annual re-selection date, separately for each stock \\
\bottomrule
\end{tabular}

\vspace{0.4em}
\begin{minipage}{0.97\textwidth}
\footnotesize
\textit{Notes:} Lookback is the number of past daily observations in the LSTM input sequence. Where a study reports several values, the table gives those used in its main results.
\end{minipage}
\end{table}
. Requires booktabs.
% Compiled 29 Oct 2026 from the papers' full texts; checked again 11 Nov 2026.
\begin{table}[!htbp]
\centering
\caption{Lookback lengths in LSTM volatility forecasting studies}
\label{tab:lookbacks}
\footnotesize
\setlength{\tabcolsep}{4pt}
\begin{tabular}{@{}p{3.6cm}p{3.9cm}p{2.0cm}p{5.6cm}@{}}
\toprule
Study & Data & Lookback (days) & How the lookback was chosen \\
\midrule
\textcite{petrozziello-2022} & 28 DJIA stocks and SPY; NASDAQ-100 stocks & 20 & Grid search over 5--100 on the first year of data; a single value reported \\
\textcite{chen-2022} & CSI 300 and S\&P 500 index futures & 20 & ``After some initial tests''; tests not reported \\
\textcite{rosenbaum-2022} & 862 US stocks, pooled & 22 & Stated, ``to make it comparable with the HAR model'' \\
\textcite{zhang-2023} & US stocks, intraday RV, pooled & 21 & Stated; the same 21-day window for every model \\
\textcite{roszyk-2024} & S\&P 500 index & 22 & Stated, ``the approximate number of trading days in a month'' \\
\textcite{perekhodko-2025} & S\&P 500 index & 21 & Stated; 5 and 42 checked afterwards \\
\textcite{souto-2023} & S\&P 500 index & 21, 63 or 252 & Random search over 21--252 on a validation split; differs by model \\
\textcite{tang-2024} & 213 cryptocurrencies, pooled & 7 or 30 & Cross-validation within the training sample \\
\textcite{rodikov-2022} & Stocks, S\&P 500, Bitcoin, Ether & 1--50 & Validation search; selected values not reported \\
\textcite{bucci-2020} & S\&P 500, monthly RV & Not reported & Information criteria for lags of RV in the predictor set; sequence length of the recurrent networks not stated \\
\midrule
This paper & 27 DJIA stocks & 6--63 & AIC on all data before each annual re-selection date, separately for each stock \\
\bottomrule
\end{tabular}

\vspace{0.4em}
\begin{minipage}{0.97\textwidth}
\footnotesize
\textit{Notes:} Lookback is the number of past daily observations in the LSTM input sequence. Where a study reports several values, the table gives those used in its main results.
\end{minipage}
\end{table}
